%% paper/sections/pre_registration.tex
%% Draft only. Every number traced via inline LaTeX comment to its source
%% file. Covers what was planned as "Experimental Setup" (datasets, protocol,
%% conditions) plus the pre-registration discipline itself.
\label{sec:pre-registration}

\subsection{Datasets and Protocol}
\textbf{MELD}~\cite{poria2019meld}, multi-party television dialogue,
%% outputs/meld_phase1_2_summary/REPORT.md: cache/MELD_MANIFEST.json 13706 clips (train 9988, dev 1108, test 2610)
13706 clips (train 9988 / dev 1108 / test 2610), 3-class sentiment. Primary
entity budget $E{=}6$ (cached at $E_{\max}{=}8$), $S{=}4$. Identity recovery
is agglomerative face-embedding clustering (shot cuts break IoU tracking),
threshold tuned on dev only, achieving
%% outputs/meld_phase1_2_summary/REPORT.md: "Best threshold 0.55 (purity/accuracy against the character references: 62.2%..."
62.2\% clustering purity against a 6-character reference set (16.7\% chance
floor) --- imperfect identity recovery is a limitation discussed in
\S\ref{sec:limitations}, not a claim of clean tracking. MELD passed a dataset
admission procedure (\S\ref{sec:dataset-audit}) before being promoted to
primary: cross-split near-duplicate rate at 0.95 cosine similarity effectively
zero, trivial-feature-probe macro-F1
%% outputs/dataset_admission/meld/admission_report.json: trivial_probe.val_macro_f1
0.4155.

\textbf{DAiSEE}~\cite{gupta2016daisee}, single-subject e-learning webcam
video,
%% logs/GATES.md: DAiSEE 8571/8571 clips (train 5358 / val 1429 / test 1784)
8571 clips (train 5358 / val 1429 / test 1784), 4-class Engagement. Primary
entity budget locked at $E{=}1$: the detector's raw output contains multiple
candidate person-tracks per clip (background artifacts, brief ID switches),
but only the single top-confidence$\times$coverage track corresponds to the
annotated subject, so $E{=}1$ discards noise rather than signal. With one
entity slot there is no second slot to shuffle, so A2 and H1 (persistence)
are not defined for this dataset (\S\ref{sec:results}). DAiSEE's own
admission-test trivial-probe score,
%% outputs/dataset_admission/daisee/admission_report.json: trivial_probe.val_macro_f1
macro-F1 0.2177 on its subject-disjoint split, is the reference floor used to
validate the admission procedure itself (\S\ref{sec:dataset-audit}).

\textbf{OUC-CGE}~\cite{lu2025ouccge} was the originally intended primary
dataset and failed the same admission procedure; \S\ref{sec:dataset-audit}
reports the audit in full. It is retained in this paper as a rejected-dataset
case study, not as a results-table entry.

\subsection{Conditions}
A0--A5 and mask-only share the EPT-Former (\S\ref{sec:method}) or a matched
simplification of it, differing only in what tokens they receive: \textbf{A0}
--- fixed spatial grid, no entities, matched token budget; \textbf{A1} ---
full EPT, entity-localized, persistent identity (primary condition);
\textbf{A2} --- EPT with entity identity shuffled per segment
(\S\ref{sec:shuffle-control}); \textbf{A3} --- temporal attention only
(social attention disabled); \textbf{A4} --- social attention only (temporal
attention disabled); \textbf{A5} --- mean-pooled entity tokens through an
MLP, no attention at all; \textbf{mask-only} --- the presence mask alone,
with no visual features, testing whether who-is-onscreen leaks label
information by itself. A \textbf{trivial-feature probe} (linear classifier
on a single mean-pooled clip embedding, no temporal/entity/attention
structure) is reported as a standing floor in every results table, not as a
condition that went through recipe search.

\subsection{Pre-registration}
Hypotheses, decision rules, and effect-size floors were committed to
\texttt{\seqsplit{docs/DECISION\_RULES.md}} (MELD) and
\texttt{\seqsplit{docs/DECISION\_RULES\_DAISEE.md}} (DAiSEE) before any
test-split evaluation, with amendments appended and dated, never edited in
place. This is stated here as a credibility mechanism, not a formality:
every number in \S\ref{sec:results} that bears on a hypothesis was computed
after its decision criterion was frozen, not the reverse.

\textbf{H1 (persistence, MELD only).} Not defined for DAiSEE ($E{=}1$,
\S\ref{sec:method}).
\[
A1 - A2 \;\geq\; \text{EFFECT\_FLOOR} \quad \text{(macro-F1)}
\]

\textbf{H2 (entity structure).} MELD uses macro-F1 directly; DAiSEE
restricts macro-F1 to classes 1--3 (see below).
\[
A1 - A0 \;\geq\; \text{EFFECT\_FLOOR} \quad \text{(macro-F1)}
\]

\textbf{Decision rule (MELD, three branches).} A branch fires only if both
the point estimate clears \texttt{EFFECT\_FLOOR} \emph{and} a paired
bootstrap over test items (10{,}000 resamples) gives a 95\% CI excluding
zero.
\begin{align*}
\text{(a)}&\quad A1 - A2 \geq \text{EFFECT\_FLOOR} \\
\text{(b)}&\quad |A1{-}A2| < \text{EFFECT\_FLOOR} \ \text{and H2 supported} \\
\text{(c)}&\quad A1 < A2 - \text{EFFECT\_FLOOR}
\end{align*}
Branch (a): persistence is the mechanism. Branch (b): entity localization
and token budget matter, persistence does not. Branch (c): hypothesis
refuted.

\textbf{Decision rule (DAiSEE, H2 only, two conditions).} Same paired
requirement as above.
\begin{align*}
\text{(a)}&\quad \text{margin} \geq \text{EFFECT\_FLOOR} \\
\text{(b)}&\quad |\text{margin}| < \text{EFFECT\_FLOOR} \\
\text{(c)}&\quad \text{margin} < -\text{EFFECT\_FLOOR}
\end{align*}
Branch (a): supported. Branch (b): null/underpowered. Branch (c): reversed.

\textbf{Effect floors, computed blind.} For each dataset, seed-level and
item-level macro-F1 variance was bootstrapped on dev/val \emph{before}
inspecting which condition scored higher --- only dispersion was ever
recorded --- then \texttt{EFFECT\_FLOOR} was set to twice the pooled SE,
floored at $0.02$ and rounded up:
\[
\text{EFFECT\_FLOOR} = \max(0.02,\ 2\,\text{SE})
\]
MELD's pooled SE was $0.0170$, giving $\text{EFFECT\_FLOOR}=0.04$
%% docs/DECISION_RULES.md 2026-08-15 amendment ("Effect floor computed"): pooled_SE=0.0170
%% commit a6d19a8 (procedure fixed) / 8c156a9 (floor computed)
(\texttt{a6d19a8}/\texttt{8c156a9}). DAiSEE's pooled SE was $0.0138$, giving
$\text{EFFECT\_FLOOR}=0.03$
%% docs/DECISION_RULES_DAISEE.md 2026-08-15 amendment: pooled_SE=0.0138
%% commit 00dcc1d (pre-registration) / 7c2c0d0 (floor computed)
(\texttt{00dcc1d}/\texttt{7c2c0d0}). Both studies are, by this procedure,
explicitly underpowered for the originally-hoped 0.02 margin --- stated in
the amendment before either floor was used to evaluate anything, not
discovered after the fact.

\textbf{DAiSEE primary-metric amendment.}
%% commit 347115b, committed before Phase 4 ran
Committed before DAiSEE's test evaluation (\texttt{347115b}), the primary
metric for DAiSEE's H2 was restricted to macro-F1 over classes 1--3, excluding
class 0 from the average (reported separately as a raw count and raw F1 in
every table, never dropped silently). Class 0 has
%% docs/DECISION_RULES_DAISEE.md: test distribution {0:4,...}
4 test clips --- too few to support a macro-averaged claim --- and scored
%% docs/LOCKED_RECIPE_DAISEE.md: class 0 F1 == 0.000 in 9 of 16 grid runs, val
0.000 F1 in 9 of 16 dev-search grid points despite 23 val clips to calibrate
against. \texttt{EFFECT\_FLOOR} was explicitly \emph{not} recomputed for the
new metric, to avoid letting the threshold move with information gathered
after the original blind bootstrap. The amendment used only val-split
behavior and test-set item counts; no test-split performance had been
observed at the time it was written.
